% CMS coordinate system, west perspective. Adapted from Izaak Neutelings,
% https://tikz.net/axis3d_cms/ (CC BY-SA 4.0, (c) 2021 TikZ.net): the geometry is his,
% the colours are Paul Tol's "vibrant" scheme and the fonts those of the paper (newtx).
\documentclass[border=3pt,tikz]{standalone}
\usepackage{newtxtext,newtxmath}
\usepackage{physics}
\usepackage{tikz}
\usepackage{tikz-3dplot}
\usetikzlibrary{arrows.meta}
\usetikzlibrary{3d}
\usepackage{xcolor}
\usetikzlibrary{bending}
\usetikzlibrary{angles,quotes}

% Paul Tol, vibrant
\definecolor{tolblue}{HTML}{0077BB}
\definecolor{tolcyan}{HTML}{33BBEE}
\definecolor{tolteal}{HTML}{009988}
\definecolor{tolorange}{HTML}{EE7733}
\definecolor{tolred}{HTML}{CC3311}
\definecolor{tolmagenta}{HTML}{EE3377}
\definecolor{tolgrey}{HTML}{BBBBBB}

% The figure's roles; change these four lines to recolour the whole figure.
\colorlet{detcol}{tolblue}   % detector outline, end faces and the eta labels
\colorlet{veccol}{tolred}    % the constituent momentum and its projections
\colorlet{pipecol}{black!50} % beam pipe
\colorlet{geocol}{black!50}  % compass and LHC labels
\colorlet{jetcol}{black}    % the jet on the barrel: cone, footprint, constituents

\colorlet{mycontour}{detcol!3}
\colorlet{mycontour2}{detcol!6}
\tikzset{>=latex}
\tikzstyle{axis}=[->,thick,line cap=round]
\tikzstyle{detector}=[thick,draw=detcol,rotate around z=\ang]
\tikzstyle{beam pipe}=[draw=pipecol,fill=black!10,rotate around z=\ang]
\tikzstyle{detector surface}=[detcol!12,opacity=0.5,rotate around z=\ang]
\renewcommand*{\vv}[1]{\vec{\mkern0mu#1}}

\begin{document}
\tdplotsetmaincoords{75}{50}
\begin{tikzpicture}[scale=2.8,tdplot_main_coords,rotate around x=90]

  % VARIABLES
  \def\rvec{\L/2/cos(\thetavec)}
  \def\thetavec{18}
  \def\phivec{60}
  \def\L{3.3}    % detector length
  \def\R{0.75}   % detector cylinder radius
  \def\l{4.3}    % beam pipe length
  \def\r{0.04}   % beam pipe radius
  \def\rt{0.042} % beam pipe radius + line thickness
  \def\xmax{1}   % maximum x axis
  \def\ymax{1}   % maximum y axis
  \def\zmin{-\l/2-0.2} % minimum z axis
  \def\zmax{\l/2+0.3}  % maximum z axis
  \def\w{0.3}
  \coordinate (O) at (0,0,0);
  \coordinate (Z) at (0,0,\L/2);
  \tdplotsetcoord{O'}{0.022}{\thetavec}{\phivec}
  \tdplotsetcoord{O''}{0.018}{90}{\phivec}
  \tdplotsetcoord{P}{\rvec}{\thetavec}{\phivec}

  % CYLINDER behind
  \def\ang{18}
  \fill[detcol!5,rotate around z=\ang]
    (0,\R,\L/2) --++ (0,0,-\L) arc(90:270:\R) --++ (0,0,\L) arc(270:90:\R) -- cycle;
  \fill[detector surface]
    (0,0,\L/2) --++ (0,\R,0) arc(90:270:\R) -- cycle;
  \fill[detector surface]
    (0,0,-\L/2) --++ (0,\R,0) arc(90:270:\R) -- cycle;
  \tdplotdrawarc[detector]{(0,0,\L/2)}{\R}{0}{360}{}{}
  \tdplotdrawarc[detector,thin]{(0,0,-\L/2)}{\R}{0}{360}{}{}
  \draw[detector,thin]
    (90:\R) arc (90:270:\R);
  \draw[detector,very thin]
    (90:\R) arc (90:90-\ang:\R);
  \draw[detector] (0,0,-\L/2)++(90:\R) --++ (0,0,\L);
  \draw[detector] (0,0,-\L/2)++(-90:\R) --++ (0,0,\L);

  % BEAM PIPE
  \tdplotdrawarc[beam pipe]{(0,0,\l/2)}{\r}{0}{360}{}{}
  \draw[beam pipe]
    (0,\r,-\l/2) -- (0,\r,-0.2*\l) -- (90:0.5*\r)
    -- (0,\r,0.2*\l) -- (0,\r,0.5*\l) arc(90:-90:\r)
    -- (0,-\r,0.2*\l) -- (-90:0.5*\r) --
    (0,-\r,-0.2*\l) -- (0,-\r,-\l/2) arc(-90:90:\r);
  \draw[beam pipe] (0,0,\l/2) circle(\r);

  % AXES
  \draw[axis,-] (0,0,\zmin) -- (0,0,0);
  \fill[black] (O) circle(0.5pt) node[right=1,below=1] {IP};
  \draw[axis] (0,0,0.020) -- (0,0,\zmax) node[right=3,below=1]{$z$};
  \draw[axis] (0,0.019,0) -- (0,\ymax,0) node[below left]{$y$};
  \draw[dashed,veccol] (O'')  -- (Pxy);
  \draw[axis] (0.022,0,0) -- (\xmax,0,0) node[below=1,right=-2]{$x$};

  % LABELS
  \node[detcol,above] at (0,\ymax,0) {$\eta=0$};
  \node[detcol,above=3] at (0,\R,0.3*\L) {$\eta>0$};
  \node[detcol,above=3] at (0,\R,-0.2*\L) {$\eta<0$};
  \node[detcol,below=1,left] at (0,0,\zmax) {$\eta=+\infty$};
  \node[detcol,above=1,right] at (0,0,\zmin) {$\eta=-\infty$};

  % JET: the chosen jet projected on the barrel (cone from the IP, footprint, constituents)
  \input{@POINTS@}
  % VECTORS
  \draw[dashed,veccol] (Px) -- (Pxy);
  \draw[->,veccol,line cap=round] (O'') -- (Pxy) node[anchor=170,inner sep=2pt] {$\vv{p}_\mathrm{T}$};
  \draw[dashed,veccol,line join=round] (Py) -- (Pxy) -- (P);
  \draw[->,veccol,line cap=round] (O') -- (P);

  % CYLINDER front
  \draw[beam pipe,fill=none] (0,\r,-\l/2) arc(90:-90:\r);
  \fill[detector surface]
    (0,\rt,\L/2) --++ (0,\R-\rt,0) arc(90:-90:\R) --++ (0,\R-\rt,0) arc(-90:90:\rt);
  \fill[detector surface]
    (0,\rt,-\L/2) --++ (0,\R-\rt,0) arc(90:-90:\R) --++ (0,\R-\rt,0) arc(-90:90:\rt);
  \tdplotdrawarc[detector]{(0,0,\L/2)}{\R}{-90}{90}{}{}
  \tdplotdrawarc[detector]{(0,0,-\L/2)}{\R}{-90}{90}{}{}
  \draw[beam pipe,fill=none] (0,\r,\l/2) arc(90:-90:\r);
  \draw[detector,very thin]
    (90-\ang:\R) arc (90-\ang:-90:\R);
  \node[veccol,anchor=-30,inner sep=2pt] at (P) {$\vv{p}$};
  \draw[dashed,veccol] (Pz) -- (P);

  % ANGLES
  \tdplotdrawarc[thick,mycontour]
    {(O)}{0.2}{4}{0.7*\phivec}{}{}
  \tdplotdrawarc[->]{(O)}{0.2}{0}{\phivec}
    {anchor=-150,inner sep=1pt}{$\phi$}
  \tdplotdrawarc[->,rotate around z=\phivec-90,rotate around y=-90]
    {(O)}{0.88}{0}{\thetavec}
    {anchor=-20,inner sep=2pt}{$\theta$}
  \tdplotdrawarc[thick,detcol!4,rotate around z=\phivec-90,rotate around y=-90]
    {(O)}{0.3}{88}{0.5*(90+\thetavec)}{}{}
  \tdplotdrawarc[-{>[flex'=1]},rotate around z=\phivec-90,rotate around y=-90,line cap=round]
    {(O)}{0.3}{90}{\thetavec}
    {anchor=-40,inner sep=1pt}{$\eta$}
  \draw[detcol] (0,0,\L/2) --++ (\R,0,0);
  \tdplotdrawarc[thick,mycontour2]
    {(Z)}{0.2}{4}{0.7*\phivec}{}{}
  \tdplotdrawarc[->]{(Z)}{0.2}{0}{\phivec}
    {anchor=-150,inner sep=1pt}{$\phi$}

  % COMPASS and LHC (CMS geography; delete in Illustrator if unwanted)
  \begin{scope}[shift={(1.1*\R,-\R,0.2*\L)},rotate around y=12]
    \draw[<->,geocol] (-\w,0,0) -- (\w,0,0);
    \draw[<->,geocol] (0,0,-\w) -- (0,0,\w);
    \node[above=1,left,geocol,scale=0.6] at (-\w,0,0) {N};
    \node[below=3,left=-2,geocol,scale=0.6] at (0,0,\w) {Jura};
    \node[below=1,right,geocol,scale=0.6] at (\w,0,0) {ATLAS};
  \end{scope}
  \draw[->,thick,geocol] (1.4*\w,-\R,-0.1*\L) --++ (2*\w,0,0)
    node[below=1,right,scale=0.8,align=center] {center of\\[-1pt]the LHC};

\end{tikzpicture}
\end{document}
